\documentclass[10pt]{article}
\usepackage[letterpaper,margin=0.82in]{geometry}
\usepackage[T1]{fontenc}
\usepackage[utf8]{inputenc}
\usepackage{lmodern}
\usepackage{microtype}
\usepackage{amsmath,amssymb}
\usepackage{booktabs,longtable,array}
\usepackage{calc}
\usepackage{xurl}
\usepackage[hidelinks]{hyperref}
\usepackage{newunicodechar}
\providecommand{\tightlist}{\setlength{\itemsep}{0pt}\setlength{\parskip}{0pt}}
\providecommand{\real}[1]{#1}
\newunicodechar{≡}{\ensuremath{\equiv}}
\newunicodechar{≈}{\ensuremath{\approx}}
\newunicodechar{⊆}{\ensuremath{\subseteq}}
\newunicodechar{Σ}{\ensuremath{\Sigma}}
\newunicodechar{Γ}{\ensuremath{\Gamma}}
\newunicodechar{ρ}{\ensuremath{\rho}}
\newunicodechar{×}{\ensuremath{\times}}
\newunicodechar{→}{\ensuremath{\to}}
\newunicodechar{∧}{\ensuremath{\land}}
\newunicodechar{≃}{\ensuremath{\simeq}}
\newunicodechar{∞}{\ensuremath{\infty}}
\newunicodechar{∈}{\ensuremath{\in}}
\newunicodechar{·}{\ensuremath{\cdot}}
\newunicodechar{−}{\ensuremath{-}}
\newunicodechar{≥}{\ensuremath{\ge}}
\newunicodechar{≤}{\ensuremath{\le}}
\newunicodechar{∘}{\ensuremath{\circ}}
\newunicodechar{□}{\ensuremath{\square}}
\title{Budgets as Grades: Which Equal-Meaning Inputs a Budgeted Reasoner Distinguishes}
\author{Jinu Jang\\
Hankuk University of Foreign Studies\\
\texttt{jinu0633@hufs.ac.kr}}
\date{September 2026}
\begin{document}
\maketitle
\begin{abstract}
A recognizer that answers Horn entailment queries with a fixed computation budget does not treat logically identical inputs alike. We ask \emph{which} of them it separates, and answer with a graded object. The closure operator of propositional Horn logic factors into stages \(T_k\), the atoms derivable in \(k\) parallel rounds of forward chaining, with \(T_0=\mathrm{id}\), \(T_j\circ T_k=T_{j+k}\), and limit the closure. We prove in Lean 4, without Mathlib, that the limit is sound and complete, that budget-\(k\) behavioural identity is blind to premise order and repetition at every grade, and that a trace is separated from a redundant extension at budget \(k\) exactly when the extension moves some derivation across the budget. This yields a preregistered prediction: a derivable clause that shortens the target derivation changes a budgeted recognizer's decision at a low budget and not at a high one, while an equally redundant clause that preserves depth changes nothing at any budget. On 200 fresh cases the prediction held for a learned iterative reasoner (interaction \(0.885\), 95\% CI \([0.84,0.93]\)) whose decisions coincided with the \(k\)-round symbolic reasoner at every budget. A second, independent prediction, the composition law read on inputs (``a hint buys exactly its depth''), was then tested on the same weights and failed at tight budgets (agreement \(0.90\) and \(0.83\) against a preregistered \(0.95\) criterion), while holding for a model trained under a smaller budget; every violation was a lost derivation, so the learned operator satisfies the lax inequality \(T_k\circ T_j\subseteq T_{j+k}\) but not the equation. Two language models (0.5B and 1.5B parameters) did not read the table at the decision level; at the margin level their sensitivity to redundant clauses was dominated by lexical overlap with the query, with a small residual in the depth direction. The graded structure therefore describes what a budgeted reasoner distinguishes; whether it also obeys the structure's composition law is a property of training rather than architecture, and was decided by testing the law rather than by assuming it.
\end{abstract}

\section{Introduction}\label{introduction}

Two inputs with the same meaning are, to a recognizer, two inputs. The recognition-paths framework \cite{rp2026} makes this precise for propositional Horn logic. A \emph{recognizer} is any map from an ordered premise trace and a query to an observation. It induces a behavioural identity \(\approx_\rho\) on traces (no context and no query separates them); logical identity \(\equiv_L\) is theory equivalence (the same consequences for every query); both are congruences on the free monoid of traces, and the recognizer \emph{respects logic} exactly when \(\equiv_L\) refines \(\approx_\rho\), equivalently when a unique monoid morphism from the logical quotient to the behavioural quotient exists (the Recognition Factorization Theorem, machine-checked in Lean).

The empirical half of that programme \cite{ppi2026} found that no recognizer examined respects logic: language models up to 1.5B parameters and small constructed models alike distinguish premise permutations, repetitions, and, most robustly, a trace from its extension by a \emph{derivable} clause, a clause that changes no consequence. A set encoder that is permutation-invariant by architecture is exactly the syntactic quotient and only partly the semantic one; an equivalence objective does not change this. That is a negative result about identification.

This paper asks the positive question it leaves open. If a recognizer does not identify all logically identical inputs, \emph{which} of them does it separate, and is there a mathematical object that predicts the answer rather than describes it after the fact? The programme's stated rule is that the object must come after an observed law, not before: observation, then equivalence, then quotient, then operations and equations, then an algebraic theory, and only then a monad. The results of \cite{ppi2026} gave the first three. This paper supplies an object that generates equations, tests two of them under preregistration, and reports that one held and one failed.

The object is the \emph{graded closure}. Forward chaining, stopped after \(k\) parallel rounds, is an operator \(T_k\) on sets of atoms. The family \((T_k)_{k\in\mathbb N}\) has \(T_0=\mathrm{id}\) and \(T_j\circ T_k=T_{j+k}\); it is an \(\mathbb N\)-graded monad on the poset of atom sets \cite{smirnov2008,fujii2016} whose limit \(T_\infty\) is the closure operator of Horn logic. Only the limit is idempotent. Reading the grade as a \emph{computation budget} turns the algebra into predictions about recognizers with a budget knob.

\textbf{Contributions.}

\begin{enumerate}
\def\labelenumi{\arabic{enumi}.}
\tightlist
\item
  \emph{Theory.} A Lean 4 development (core library only) of the graded closure: the graded-monad laws, soundness and completeness of the limit, a budget-\(k\) behavioural identity \(\approx_k\) that is blind to order and repetition at every grade, the theorem that a trace and its extension by one clause are separated at budget \(k\) if and only if the extension moves some query's derivation from depth \(>k\) to depth \(\le k\), and the composition law read on inputs: pre-saturating the hypotheses by \(j\) rounds and reading with budget \(k\) equals reading with budget \(j+k\).
\item
  \emph{A held prediction.} The separation theorem, preregistered as an interaction between two budgets on a frozen table of 200 fresh Horn cases, held for a learned iterative reasoner with a rounds budget, with the full predicted shape and with exact agreement with the symbolic \(k\)-round reasoner at every budget.
\item
  \emph{A failed prediction.} The composition law, preregistered on the same weights, failed at the two tight budget pairs and held at the two others; all violations were lost derivations from multi-atom hypothesis sets. A model trained under a smaller budget satisfied the law. The learned reasoner is exactly graded in what it separates and only laxly graded in how it composes.
\item
  \emph{Language models.} At 0.5B and 1.5B parameters, single forward pass, the models do not read the table at the decision level; on the margin level the variable that predicts their sensitivity to redundant clauses is lexical overlap with the query, not depth.
\end{enumerate}

Everything is frozen (SHA-256 locks on tables, commit hashes on code, weights on Hugging Face with training logs) and every learned-model result was preceded by a committed preregistration naming the comparison, the criterion, the controls and the sentence to be reported on failure.

\section{Related work}\label{related-work}

\emph{Order sensitivity of language models.} Chen et al.~\cite{chen2024premise} showed that premise order changes accuracy on deductive tasks even when the logic is unchanged; metamorphic and reformulation testing \cite{zhou2026lgmt,gu2026crtbench} and parameterized logical benchmarks \cite{essebbani2026robustness} measure the same phenomenon through accuracy differences. Our instrument reads order-\emph{relations} on the answer logits so that identity is exact rather than approximate, and compares language models and constructed recognizers on the same frozen table \cite{ppi2026}. Set-LLM \cite{egressy2025setllm} and order-centric augmentation \cite{he2025order} aim at permutation invariance; the companion paper showed that permutation invariance is strictly weaker than logical invariance, and the present one shows that the gap has a graded structure.

\emph{Depth and computation.} PrOntoQA \cite{saparov2023prontoqa} and related synthetic suites vary proof depth and find accuracy falling with depth; Merrill and Sabharwal \cite{merrill2024cot} bound what transformers compute per step and what chain of thought adds; Barceló et al.~\cite{barcelo2020gnn} characterise message-passing networks by the logic of \(k\)-hop neighbourhoods. Our reasoner is a message-passing network by design, so the one-hop-per-round bound is part of its specification, and we say so in advance. What is tested is whether training realises the graded reader rather than a shortcut, whether the gap closes exactly at the predicted budget, and whether the \emph{composition} law, which no locality bound implies, holds.

\emph{Graded monads and closure operators.} Graded (parametric effect) monads \cite{smirnov2008,fujii2016,katsumata2014parametric} index a monad by a monoid; the closure operator of Horn logic is the classical least-model construction \cite{dowling1984horn}, and its truncation at \(k\) rounds is the standard bounded forward chaining. The novelty is not the algebra, which is elementary, but its use as a source of \emph{falsifiable equations} about learned recognizers, and the observation that one equation is a property of architecture (identification) while the other is a property of training (composition).

\emph{Automata and identification.} Behavioural identity as Nerode equivalence and finite-test identification as closure \cite{nerode1958,angluin1987learning} underlie the companion framework; graded identity \(\approx_k\) is a family of such equivalences indexed by budget, coarser than the syntactic quotient and finer than \(\equiv_L\) for every finite \(k\).

\section{Setting}\label{setting}

We recall the definitions of \cite{rp2026} that are used here. Fix a type of atoms. A Horn clause is \(a_1\wedge\dots\wedge a_m\to b_1\wedge\dots\wedge b_n\) with atoms on both sides; a \emph{trace} \(w\) is a finite list of clauses; its \emph{theory} \(\Gamma(w)\) is the set of clauses occurring in it. A \emph{query} is a pair \((H\Rightarrow g)\) of a list of hypothesis atoms and a goal atom. Semantic entailment \(\Gamma\models(H\Rightarrow g)\) holds when every model of \(\Gamma\) that makes \(H\) true makes \(g\) true. \emph{Logical identity} is \(u\equiv_L w\) iff \(\Gamma(u)\models q\leftrightarrow\Gamma(w)\models q\) for every query \(q\); it is a congruence for concatenation, and permuting, repeating, or appending a derivable clause to a trace preserves it.

A \emph{recognizer} is a map \(\rho(w,q)\) from traces and queries to an observation set. Its \emph{behavioural identity} is \(u\approx_\rho w\) iff \(\rho(xuz,q)=\rho(xwz,q)\) for all contexts \(x,z\) and queries \(q\); it is a congruence, and the quotient \(B_\rho\) is a monoid. The recognizer \emph{respects logic} when \(u\equiv_Lw\) implies \(u\approx_\rho w\), equivalently (the Recognition Factorization Theorem) when there is a unique monoid morphism \(L\to B_\rho\) commuting with the two quotient maps. The experiments of \cite{ppi2026} read language models through Boolean order relations on the answer logits (decision: YES-logit above NO-logit; preference: one query's margin above another's) so that \(\approx_\rho\) restricted to a finite test family is an exact identity of Boolean profiles, and found the refinement to fail for every recognizer tested.

\section{The graded closure}\label{the-graded-closure}

\subsection{Definitions and laws}\label{definitions-and-laws}

Let \(S\) range over sets of atoms, represented as predicates. One parallel round of forward chaining under \(\Gamma\) is \[\mathrm{step}_\Gamma(S)=S\cup\{b : (A\to B)\in\Gamma,\ A\subseteq S,\ b\in B\},\] and the \(k\)-round operator is defined by \[T_0S=S,\qquad T_{k+1}S=\mathrm{step}_\Gamma(T_kS).\] All statements below are proved in \texttt{RecognitionPaths/Graded.lean}; Lean names are given in parentheses.

\textbf{Proposition 1 (graded-monad laws).} For all \(j,k\) and \(S\): (i) \(T_j(T_kS)=T_{j+k}S\) (\texttt{rounds\_add}); (ii) \(S\subseteq T_kS\subseteq T_{k+d}S\) (\texttt{rounds\_extensive}, \texttt{rounds\_le\_add}); (iii) \(T_k\) is monotone in \(S\) and in \(\Gamma\) (\texttt{rounds\_mono}, \texttt{rounds\_mono\_theory}).

\emph{Proof.} (i) by induction on \(j\), using \(T_{j+1}=\mathrm{step}\circ T_j\) and \((j+1)+k=(j+k)+1\); (ii) and (iii) by induction on \(k\), since \(\mathrm{step}\) is extensive and monotone in both arguments. \(\square\)

Thus \((T_k)\) is an \(\mathbb N\)-graded monad on the poset of atom sets, with unit the identity at grade \(0\) and multiplication the identity \(T_jT_k=T_{j+k}\); it is \emph{not} a monad at any single positive grade, since \(T_kT_k=T_{2k}\neq T_k\) in general.

\textbf{Theorem 2 (the limit is the closure).} Let \(T_\infty S=\bigcup_kT_kS\). (i) \emph{Soundness:} every atom of \(T_kS\) holds in every model of \(\Gamma\) that contains \(S\) (\texttt{rounds\_sound}). (ii) \emph{Completeness:} \(T_\infty S\) is a model of \(\Gamma\) (\texttt{limit\_models}); hence \(\Gamma\models(H\Rightarrow g)\iff\exists k,\ g\in T_kH\) (\texttt{entails\_iff\_exists\_rounds}).

\emph{Proof.} (i) by induction on \(k\): a clause whose body holds in the model has its head in the model. (ii) For a clause \(A\to B\) of \(\Gamma\) with \(A\subseteq T_\infty S\), each atom of the finite list \(A\) lies in some \(T_{k_a}S\); by nestedness all lie in \(T_KS\) for \(K=\sum k_a\), so \(B\subseteq
T_{K+1}S\). Completeness follows by instantiating the definition of \(\models\) at the model \(T_\infty H\), which contains \(H\) at grade \(0\). \(\square\)

\subsection{Graded identity}\label{graded-identity}

Write \(\Gamma\vdash_k(H\Rightarrow g)\) for \(g\in T_kH\) (\texttt{EntailsK}). \emph{Budget-\(k\) identity} of traces is \[u\approx_kw\iff\forall q,\ \Gamma(u)\vdash_kq\leftrightarrow\Gamma(w)\vdash_kq\qquad(\texttt{GradedEquiv}).\]

\textbf{Proposition 3.} For every \(k\): (i) \(\approx_k\) is an equivalence relation; (ii) permuting or repeating a trace preserves it (in Lean, \texttt{gradedEquiv\_of\_perm} and \texttt{gradedEquiv\_dup}); (iii) if \(u\approx_kw\) for every \(k\) then \(u\equiv_Lw\) (\texttt{logicalEquiv\_of\_gradedEquiv}), and conversely \(u\equiv_Lw\) iff for every query the sets of grades at which \(u\) and \(w\) answer it are both empty or both nonempty (\texttt{logicalEquiv\_iff\_limit}).

\emph{Proof.} (ii) \(T_k\) depends on the trace only through its theory, and \(\Gamma(u)=\Gamma(w)\) for permutations and repetitions. (iii) from Theorem 2(ii). \(\square\)

The converse of (iii) fails at every fixed grade, and the failure is exactly characterised:

\textbf{Theorem 4 (separation of an extension).} For a trace \(w\), a clause \(c\) and a budget \(k\), \[w\approx_kw{+}[c]\iff\forall q:\ \Gamma(w{+}[c])\vdash_kq\ \Rightarrow\ \Gamma(w)\vdash_kq\] (\texttt{gradedEquiv\_append\_iff}). Equivalently, \(w\) and \(w{+}[c]\) are separated at budget \(k\) iff some query is answered within \(k\) rounds after the extension and not before it (\texttt{not\_gradedEquiv\_append\_iff}).

\emph{Proof.} Adding a clause can only add derivations within a fixed budget (monotonicity in \(\Gamma\)), so one implication of \(\approx_k\) always holds and the other is the stated condition. \(\square\)

\textbf{Corollary 5 (transience).} If \(c\) is derivable from \(w\), then for every query \(q\) there is a grade \(K\) with \(\Gamma(w{+}[c])\vdash_{K+d}q\leftrightarrow\Gamma(w)\vdash_{K+d}q\) for all \(d\) (\texttt{gradedEquiv\_append\_derivable\_eventually}).

\emph{Proof.} If \(\Gamma(w)\models q\), take \(K\) from completeness; otherwise both sides are false at every grade by soundness and \(w{+}[c]\equiv_Lw\). \(\square\)

Theorem 4 says what the budgeted recognizer distinguishes: a redundant extension is visible exactly when it \emph{shortens a derivation across the budget}. Corollary 5 says the visibility disappears once the budget covers the depth. Read as a prediction about an arbitrary recognizer with a budget knob, the theorem asserts a specific interaction between the kind of extension and the budget, and that is what is tested in Section 6.

\subsection{The budgeted recognizer and the composition law}\label{the-budgeted-recognizer-and-the-composition-law}

Let \(\rho_k(w,q)=[\Gamma(w)\vdash_kq]\) (\texttt{roundsRecognizer}). It is a theory-factoring recognizer in the sense of \cite{rp2026}, so it is invariant under permutation and repetition of any block in any context, and its behavioural identity is \(\approx_k\) in every context: \(u\approx_{\rho_k}w\iff\forall x,z,\ xuz\approx_kxwz\) (\texttt{roundsRecognizer\_contextEquiv\_iff}).

\textbf{Proposition 7 (composition law on inputs).} For all \(j,k\), hypotheses \(H\) and goal \(g\): \(g\in T_k(T_jH)\iff g\in T_{j+k}H\) (\texttt{entailsK\_presaturate}). On the recognizer, \[\rho_k(w,\ T_jH\Rightarrow g)=\rho_{j+k}(w,\ H\Rightarrow g):\] pre-saturating the hypotheses by \(j\) rounds and reading with budget \(k\) is reading with budget \(j+k\). A hint buys exactly its depth.

Proposition 7 is a restatement of Proposition 1(i), but as a statement about inputs it is testable on any recognizer that accepts a hypothesis \emph{set} and has a budget knob, and it is not implied by the locality bound that makes Theorem 4's pattern available to a message-passing architecture: a model may derive correctly from single atoms at every budget and still fail to use a set of hypotheses as the algebra requires. It is therefore the second, independent prediction.

\section{Predictions and preregistration}\label{predictions-and-preregistration}

Both predictions were committed, with the rules below, before any learned recognizer was run on the corresponding table; the symbolic reasoner and the oracle were run first as deterministic control validation.

\textbf{P1 (identification).} Take a base theory \(D\) with a target query \((a\Rightarrow g)\) of depth \(d\ge3\). Let \(F\) append a derivable clause that shortens the target's depth by the least possible amount, and \(C\) append a derivable clause that preserves it; \(D\equiv_LF\equiv_LC\) by certified closure equality. For a recognizer with budget \(k\) let \(\mathrm{dis}_X(k)\) be the fraction of cases whose target decision differs between \(D\) and \(X\), and \(\Delta(k)=\mathrm{dis}_F(k)-\mathrm{dis}_C(k)\). Theorem 4 predicts \(\Delta(k)>0\) for \(k\) below the base depth and \(\Delta(k)=0\) once \(k\) covers it. The primary comparison is the interaction \(I=\Delta(2)-\Delta(4)\) for a learned reasoner trained at four rounds. \emph{Criterion:} \(I>0\) with a paired case-bootstrap 95\% confidence interval (\(B=5000\)) excluding \(0\), and \(\Delta(2)\ge0.10\). \emph{Controls:} the exact oracle and a constant recognizer (both \(\Delta=0\)), the \(k\)-round symbolic reasoner (the theorem by construction), Pythia-70M as an empirical non-reader, and a logic-change condition \(L\) (one clause deleted so the target becomes underivable) that every reading recognizer must separate.

\textbf{P3 (composition).} On the same weights, with hypothesis sets \(H_j=T_j\{a\}\) for \(j\in\{0,1,2\}\), the agreement rate between \(\rho(H_j;k)\) and \(\rho(H_0;j{+}k)\) over cases must have bootstrap lower bound \(\ge0.95\) on each of the four pairs \((j,k)\in\{(1,2),(2,1),(2,2),(1,3)\}\), chosen because the symbolic yes-rate on the target changes between budgets \(k\) and \(j+k\) there, so that a constant answer cannot pass.

Sample size: 200 logical structures; for proportion differences near \(0.15\) the paired bootstrap half-width is about \(0.05\), and for agreement near \(0.97\) about \(0.025\). Both preregistrations state the sentence to be reported on failure; the one for P3 is reported verbatim in Section 6.3.

\section{Instruments}\label{instruments}

\subsection{Tables}\label{tables}

Seven atoms \(a,\dots,g\). A base theory is five random Horn clauses (70\% single-body single-head, 15\% two-body, 15\% two-head) with target hypothesis \(a\) and goal the derivable atom of maximal depth, redrawn until the depth is at least \(3\). Candidate extensions are single-atom clauses \(x\to y\) with \(y\in T_\infty\{x\}\), not already a sub-clause of the theory. \(F\) is the sorted-first candidate among those that reduce the target depth by the least amount; \(F_1\) is the direct clause \(a\to g\) (maximal shortening and maximal lexical overlap with the query, a secondary condition); \(C\) is the sorted-first depth-preserving candidate with the same overlap with \(\{a,g\}\) as \(F\) when one exists (181 of 200 cases); \(L\) deletes the sorted-first clause whose removal makes the target underivable. Cases without an \(F\), a \(C\), an \(L\), a depth-1 atom or a non-derivable atom are dropped and the next draw used: 200 cases from 264 draws, one seed, taken in order. Realised depths: base \(3\) in 177 cases and \(4\) in 23; \(F\) reduces to \(2\) in 176, \(3\) in 23, \(1\) in 1. Four queries per case: the target \(t\), a depth-1 atom \(d_1\), a non-derivable atom \(n\), and the reversed query \(r\). Rows: \(200\times5\times4=4000\). The composition table takes the 200 base theories with \(H_0=\{a\}\), \(H_1=T_1\{a\}\) (\(|H_1|\in\{2,3,4\}\)), \(H_2=T_2\{a\}\) (\(|H_2|\in\{3,4,5\}\)), for \(t\) and \(n\): 1200 rows. Both tables carry SHA-256 locks.

\subsection{Recognizers}\label{recognizers}

\emph{Symbolic.} The exact closure oracle; the \(k\)-round reasoner \([g\in T_kH]\) for \(k\in\{0,\dots,6\}\); a constant-NO recognizer.

\emph{Iterative learned reasoner.} Each atom carries a state vector (\(d=64\)), initialised from a learned embedding plus a learned flag on the hypothesis atoms. In each round every clause computes a message from the sum and the elementwise minimum of its body states and the current head state (a two-layer MLP), messages are aggregated per head atom by an elementwise maximum (so the model is invariant to clause permutation and repetition by construction), and each atom's state is updated by a GRU cell. After \(k\) rounds the goal state and the maximum over hypothesis states are read out by an MLP into two logits. The number of rounds is the budget and can be set at evaluation. 75,202 parameters. Trained 6000 steps (batch 128, AdamW, one-cycle schedule) on random theories of two to six clauses over seven atoms with balanced labels and random atom relabelling, with every evaluation theory (\(D,F,F_1,C,L\) of every case) excluded from training up to relabelling. The primary model is trained with a budget of four rounds (in-distribution validation accuracy \(0.996\)); a secondary model with two rounds (\(0.960\)). Both were trained with single-atom hypotheses; the composition test feeds them hypothesis \emph{sets} through the same flag, out of the training format (for a single atom the modified readout is numerically identical to the trained one).

\emph{Set recognizer.} A seven-atom version of the max-pooled clause encoder of \cite{ppi2026} (168,386 parameters, validation accuracy \(0.951\)); no budget knob.

\emph{Language models.} Qwen2.5-0.5B-Instruct and Qwen2.5-1.5B-Instruct, raw prompt (no chat template), bullet rendering of the clauses with per-case atom names, one forward pass on CPU in float32, observation the pair of logits of the tokens YES and NO at the answer position; Pythia-70M (step 143000) as the non-reading control.

\subsection{Readout}\label{readout}

Decisions are Boolean (positive logit strictly above negative; ties count as NO). The extended observation is the real margin. All statistics are over the 200 cases with paired case bootstrap.

\section{Results}\label{results}

\subsection{P1 holds, and the learned reasoner is the symbolic one}\label{p1-holds-and-the-learned-reasoner-is-the-symbolic-one}

Table 1 gives the primary model at each budget.

{\def\LTcaptype{none} % do not increment counter
\begin{longtable}[]{@{}rrrrrrl@{}}
\toprule\noalign{}
budget \(k\) & acc\((D)\), target & \(\mathrm{dis}_F\) & \(\mathrm{dis}_{F_1}\) & \(\mathrm{dis}_C\) & \(\mathrm{dis}_L\) & \(\Delta\) {[}95\% CI{]} \\
\midrule\noalign{}
\endhead
\bottomrule\noalign{}
\endlastfoot
1 & 0.00 & 0.01 & 1.00 & 0.00 & 0.00 & 0.01 {[}0.00, 0.01{]} \\
\textbf{2} & 0.00 & \textbf{0.89} & 1.00 & \textbf{0.00} & 0.00 & \textbf{0.89 {[}0.84, 0.93{]}} \\
3 & 0.89 & 0.11 & 0.11 & 0.01 & 0.89 & 0.10 {[}0.06, 0.15{]} \\
\textbf{4} & 1.00 & \textbf{0.00} & 0.00 & \textbf{0.00} & 1.00 & \textbf{0.00 {[}0.00, 0.00{]}} \\
6 & 1.00 & 0.00 & 0.00 & 0.00 & 1.00 & 0.00 {[}0.00, 0.00{]} \\
\end{longtable}
}

Table 1. The four-round learned reasoner on 200 cases; decision-change rates relative to the base \(D\) on the target query. \(I=\Delta(2)-\Delta(4)=0.885\), 95\% CI \([0.84,0.93]\). \textbf{P1 holds.}

The learned reasoner's decision table coincides with that of the \(k\)-round symbolic reasoner at every budget (Table 4 in the appendix lists both). At \(k=2\) the 177 depth-3 bases are NO and their \(F\) extensions (depth 2) YES; the 23 depth-4 bases are NO on both. At \(k=3\) the depth-4 bases separate (0.11) and at \(k=4\) nothing does except the logic change. Every \(D\)--\(F\) disagreement is \(F\) correct and \(D\) wrong; the reverse never occurs. The direct clause \(F_1\) separates already at \(k=1\), because it is a depth-1 derivation, and not at \(k=4\), so lexical overlap with the query plays no role for this recognizer. The controls behave as required: oracle and constant recognizer give \(\Delta=0\); Pythia-70M gives \(\Delta=-0.02\) \([-0.04,0.01]\) and \(\mathrm{dis}_L=0.04\), i.e.~it does not read.

A caveat stated in the preregistration: information in this architecture flows one hop per round, so a \emph{correct} model must show this pattern. The empirical content is that training produced this reader rather than a shortcut (no clause-count or lexical heuristic survives at \(k=2\)), that the gap closes exactly at the predicted budget, and that the pattern is what Theorem 4 says and nothing more.

\subsection{Secondary constructed recognizers}\label{secondary-constructed-recognizers}

\emph{The two-round model.} Trained under a budget too small for its data, it shows the same interaction, \(I=0.665\) \([0.60,0.74]\), and closes the gap at \(k=4\) although it never ran four rounds in training (its rounds extrapolate; at \(k=6\) accuracy degrades to \(0.90\) and the logic change becomes partly invisible, \(\mathrm{dis}_L=0.70\)). It differs from the primary model in one informative way: at its own budget it says YES on 16\% of the depth-3 bases it cannot derive, and its decision moves on 10\% of the depth-preserving extensions \(C\). A model trained under an insufficient budget acquires a shortcut component that a merely redundant clause can trigger; depth accounts for \(0.67\) of its \(0.77\) \(F\) effect.

\emph{The set recognizer.} With no budget it identifies almost every extension at the decision level (\(\mathrm{dis}_F=0.01\), \(\mathrm{dis}_{F_1}=0.00\), \(\mathrm{dis}_C=0.00\)) and reads the logic change (\(\mathrm{dis}_L=0.94\)). Its margins nevertheless order the conditions \(F_1>F>C>0>L\) (mean shifts on the target \(+1.91\), \(+0.65\), \(+0.23\), \(-9.3\) logits, with \(F-C=0.42\) \([0.29,0.54]\)): a recognizer without rounds still orders redundant extensions by depth shortening in its extended observation.

\subsection{P3 fails for the primary model}\label{p3-fails-for-the-primary-model}

{\def\LTcaptype{none} % do not increment counter
\begin{longtable}[]{@{}lrll@{}}
\toprule\noalign{}
pair \((j,k)\) & symbolic & four-round model (primary) & two-round model \\
\midrule\noalign{}
\endhead
\bottomrule\noalign{}
\endlastfoot
\((1,2)\) & 1.000 & \textbf{0.897 {[}0.870, 0.925{]}} & 0.990 {[}0.980, 0.998{]} \\
\((2,1)\) & 1.000 & \textbf{0.830 {[}0.797, 0.863{]}} & 0.975 {[}0.960, 0.990{]} \\
\((2,2)\) & 1.000 & 0.995 {[}0.988, 1.000{]} & 0.998 {[}0.993, 1.000{]} \\
\((1,3)\) & 1.000 & 0.983 {[}0.970, 0.995{]} & 1.000 {[}1.000, 1.000{]} \\
\end{longtable}
}

Table 2. Agreement between \(\rho(H_j;k)\) and \(\rho(H_0;j{+}k)\) over 200 cases and two queries, with paired bootstrap 95\% CI.

Two of the four primary pairs fall below the criterion, so \textbf{P3 fails} for the primary model. As preregistered, we report: \emph{on 200 fresh Horn cases the learned reasoner did not satisfy the composition law \(T_k\circ T_j=T_{j+k}\) out of its training format (agreement \(0.897\) \([0.870,0.925]\) at \(j=1,k=2\) and \(0.830\) \([0.797,0.863]\) at \(j=2,k=1\)); the graded-monad description fits its budget behaviour but not its hypothesis-set behaviour, and the law is a property of the symbolic reasoner that the learned one does not inherit.}

The failure has one direction. In all 41 disagreements at \((1,2)\) and all 68 at \((2,1)\) the model answers NO from the pre-saturated hypotheses at the tight budget where the law says YES; it never answers YES where the law says NO. It reads multi-atom hypothesis sets correctly when given slack (accuracy \(1.000\) at \(k=4\) with \(H_1\) and with \(H_2\)); agreement falls with the size of the hint (at \((2,1)\): \(|H_2|=3\): \(0.73\), \(4\): \(0.63\), \(5\): \(0.55\)) and is near-perfect on the depth-4 cases (\(0.957\)), where a spare round exists. The learned operator therefore satisfies the lax inequality \(T_k\circ T_j\subseteq T_{j+k}\), with strict inclusion on 20--35\% of cases at the tight budgets: it loses part of a round when a derivation starts from a set rather than an atom. The two-round model, which had to propagate from whatever it was given within two rounds, passes all four pairs; this was not preregistered for it and is reported as an observation.

\subsection{Language models}\label{language-models}

On this rendering (seven atoms, five or six clauses) Qwen2.5-0.5B answers YES on every query of every condition (minimum margin \(1.23\)) and Qwen2.5-1.5B answers NO on every query, so the decision-level test cannot be evaluated for either; this is the ``not reading'' gate of the audit protocol in \cite{ppi2026} and is reported as not evaluable, not as a refutation. The margins carry partial reading (comparative accuracy \([m_t>m_n]\) on the base: \(0.79\) and \(0.89\)), so the extended observation can be examined, as preregistered, in an exploratory analysis.

{\def\LTcaptype{none} % do not increment counter
\begin{longtable}[]{@{}llll@{}}
\toprule\noalign{}
recognizer & \(F-C\) & \(F_1-F\) & \(L\) \\
\midrule\noalign{}
\endhead
\bottomrule\noalign{}
\endlastfoot
Pythia-70M (non-reader) & 0.010 {[}0.005, 0.014{]} & 0.004 {[}−0.001, 0.009{]} & −0.080 \\
Qwen2.5-0.5B & 0.033 {[}0.026, 0.040{]} & 0.098 {[}0.091, 0.104{]} & −0.406 \\
Qwen2.5-1.5B & 0.048 {[}0.030, 0.067{]} & 0.285 {[}0.268, 0.302{]} & −0.207 \\
set recognizer & 0.417 {[}0.289, 0.540{]} & 1.264 {[}1.124, 1.407{]} & −9.32 \\
\end{longtable}
}

Table 3. Mean signed margin shift on the target query relative to the base, in logits, over 200 cases; paired bootstrap 95\% CI. Exploratory.

For both language models the shortening clause moves the margin more than the depth-preserving clause, with intervals excluding \(0\), but by \(0.03\)--\(0.05\) logits: three to five times the non-reader's surface residual (which exists because \(F\) mentions the hypothesis atom in 198 of 200 cases and \(C\) in 131; on the 181 overlap-matched cases \(F-C\) is \(0.025\) and \(0.029\)) and an order of magnitude below the direct clause's effect (\(F_1-F\): \(0.10\) and \(0.29\)). The variable that predicts a redundant clause's visibility to these models is lexical overlap with the query; the depth direction is a small residual.

\section{Discussion}\label{discussion}

\emph{What the graded object explains.} The companion paper found that every recognizer separates a trace from its redundant extension and could not say why. Theorem 4 gives the reason for one class of recognizers: separation happens exactly when the extension moves a derivation across the budget, and it is transient in the budget. The prediction held with the full shape, including the zero on depth-preserving extensions at every budget, the zero on all extensions at sufficient budget, and the one-sided direction of every disagreement. For the learned reasoner, then, the answer to ``which equal-meaning inputs does it distinguish?'' is exact and is given by the grade.

\emph{Architecture versus training.} Two equations of the same object were tested on the same weights. The identification equation is made available by the architecture (one hop per round) and was realised by training. The composition equation is not implied by any locality bound, and the trained model violates it in a specific way: it is exactly graded in what it separates and only laxly graded in how it composes, with \(T_kT_j\subseteq T_{j+k}\) and a loss of part of a round when derivations start from sets. A model trained under a smaller budget satisfies the equation. The graded monad is therefore the right description of the \emph{quotient} the recognizer computes; whether the recognizer also carries the monad's \emph{multiplication} depends on how it was trained. In the programme's terms, the algebraic theory has been reached for one operation and refuted for one equation, by a test of the equation and not by construction.

\emph{The language-model variable.} For the two language models the graded account is not what predicts sensitivity; lexical overlap with the query is. This is consistent with the companion paper's finding that surface change outweighs a one-arrow logical change for these models, and it sharpens it: the surface variable is now named, and the depth variable is measured as a residual an order of magnitude smaller. Nothing here says that larger models or models with a genuine budget behave the same way.

\emph{The next test.} The composition law is an equation between budgets, and a single forward pass has no budget knob. A language model with chain of thought has one, the number of reasoning tokens, and Proposition 7 then makes a quantitative prediction: supplying the \(j\)-round consequences of the hypotheses as premises should be worth exactly \(j\) rounds of budget, no more. Whether the lax inequality observed here for the learned reasoner also holds for a language model, and whether it ever becomes an equation, is the natural continuation, and it needs generation rather than a single pass.

\section{Limitations}\label{limitations}

Two language models, both under two billion parameters, one rendering, no chat template, one forward pass; the decision-level LLM tests are not evaluable at this scale and the margin-level analysis is exploratory. The constructed models are small and the iterative reasoner's architecture makes the identification pattern available to a correct model, as stated in advance; the composition result for the two-round model was not preregistered. Seven atoms and five-clause bases bound the depths to \(3\) and \(4\); the budgets compared are therefore \(2\) against \(4\), and monotone behaviour across all budgets was not assumed or claimed. The Lean development is elementary and is offered to fix definitions, not as new mathematics. No claim is made that any recognizer, learned or linguistic, ``has'' a monad; the paper's claims are that a named algebraic object predicts one behaviour of a budgeted recognizer exactly, fails to predict another, and does not predict the language models examined.

\section{Artifacts and reproducibility}\label{artifacts-and-reproducibility}

Theory: \texttt{RecognitionPaths/Graded.lean} in the repository \texttt{recognition-paths} (Lean 4 v4.24.0, no Mathlib; concept DOI 10.5281/zenodo.22495808). Tables, generators, runners, analysis and preregistrations: \texttt{rq2/} and \texttt{docs/} in \texttt{proof-path-invariance} (concept DOI 10.5281/zenodo.22495706). The frozen tables have the SHA-256 locks

\begin{verbatim}
rq2_prompts.jsonl   39a82a17f9aae4027586339c58942a0114439bf0f978cf12724e754479676510
rq2b_presat.jsonl   5e42376dd5e01a7f4f6b04856a74a4596f8378ea6579786acce489928456a910
\end{verbatim}

Weights and training logs: \texttt{jinu0633/recognition-paths-recognizers}, folder \texttt{rq2/} (Hugging Face). The paper repository \texttt{graded-recognition} carries the manuscript, the exact commands, and verbatim copies of the preregistrations with their outcomes. All runs are CPU-only.

\section*{Appendix: control validation and full budget tables}\label{appendix-control-validation-and-full-budget-tables}
\addcontentsline{toc}{section}{Appendix: control validation and full budget tables}

{\def\LTcaptype{none} % do not increment counter
\begin{longtable}[]{@{}rrlrl@{}}
\toprule\noalign{}
budget \(k\) & symbolic acc\((D)\) & symbolic \(\mathrm{dis}_F\) / \(\mathrm{dis}_C\) / \(\mathrm{dis}_L\) & learned acc\((D)\) & learned \(\mathrm{dis}_F\) / \(\mathrm{dis}_C\) / \(\mathrm{dis}_L\) \\
\midrule\noalign{}
\endhead
\bottomrule\noalign{}
\endlastfoot
1 & 0.75 & 0.01 / 0.00 / 0.00 & 0.75 & 0.01 / 0.00 / 0.00 \\
2 & 0.75 & 0.89 / 0.00 / 0.00 & 0.75 & 0.89 / 0.00 / 0.00 \\
3 & 0.97 & 0.12 / 0.00 / 0.89 & 0.97 & 0.11 / 0.01 / 0.89 \\
4 & 1.00 & 0.00 / 0.00 / 1.00 & 1.00 & 0.00 / 0.00 / 1.00 \\
6 & 1.00 & 0.00 / 0.00 / 1.00 & 1.00 & 0.00 / 0.00 / 1.00 \\
\end{longtable}
}

Table 4. The \(k\)-round symbolic reasoner and the four-round learned reasoner on the same table; accuracy is over all four queries.

{\def\LTcaptype{none} % do not increment counter
\begin{longtable}[]{@{}lrrrr@{}}
\toprule\noalign{}
pair \((j,k)\) & yes-rate \((H_j;k)\) & yes-rate \((H_0;j{+}k)\) & agreement on \(t\) & agreement on \(n\) \\
\midrule\noalign{}
\endhead
\bottomrule\noalign{}
\endlastfoot
\((1,1)\) & 0.000 & 0.000 & 1.000 & 1.000 \\
\((1,2)\) & 0.685 & 0.890 & 0.795 & 1.000 \\
\((2,1)\) & 0.550 & 0.890 & 0.660 & 1.000 \\
\((2,2)\) & 0.990 & 1.000 & 0.990 & 1.000 \\
\((1,3)\) & 0.965 & 1.000 & 0.965 & 1.000 \\
\((1,4)\), \((2,3)\), \((2,4)\) & 1.000 & 1.000 & 1.000 & 1.000 \\
\end{longtable}
}

Table 5. The composition test for the four-round model by query; the symbolic reasoner agrees at 1.000 on every pair and the two-round model at 0.975 or above on every primary pair.


\begin{thebibliography}{99}
\bibitem{rp2026}
J. Jang.
Recognition paths: logical and behavioural identity for Horn traces (Lean 4 development).
Zenodo, 2026. Version 0.2.0, doi:10.5281/zenodo.22636567 (concept doi:10.5281/zenodo.22495808).

\bibitem{ppi2026}
J. Jang.
Recognition paths: when does a recognizer respect logical identity, and how would we know?
Manuscript and repository \emph{proof-path-invariance}, Zenodo, 2026. Version 0.2.0, doi:10.5281/zenodo.22636570 (concept doi:10.5281/zenodo.22495706).

\bibitem{smirnov2008}
A. L. Smirnov.
Graded monads and rings of polynomials.
\emph{Journal of Mathematical Sciences}, 151(3):3032--3051, 2008.

\bibitem{fujii2016}
S. Fujii, S. Katsumata, and P.-A. Melli\`es.
Towards a formal theory of graded monads.
In \emph{Foundations of Software Science and Computation Structures (FoSSaCS)}, LNCS 9634, pages 513--530. Springer, 2016.

\bibitem{chen2024premise}
X. Chen, R. A. Chi, X. Wang, and D. Zhou.
Premise order matters in reasoning with large language models.
In \emph{Proceedings of ICML}, 2024. arXiv:2402.08939.

\bibitem{zhou2026lgmt}
Z. Zhou, M. Li, X. Fang, X. Zhou, W. Lin, and Z. Zheng.
LGMT: Logic-grounded metamorphic testing for evaluating the reasoning reliability of LLMs.
\emph{Knowledge-Based Systems}, 348:116324, 2026. arXiv:2605.23965.

\bibitem{gu2026crtbench}
A. Gu and A. Chen.
Controlled reformulation testing for logical consistency in large language models.
arXiv:2607.14528, 2026.

\bibitem{essebbani2026robustness}
N. Es-sebbani, E. Marquer, Y. Salhi, and Z. Bouraoui.
Evaluating robustness of reasoning models on parameterized logical problems.
arXiv:2602.12665, 2026.

\bibitem{egressy2025setllm}
B. Egressy and J. Stuehmer.
Set-LLM: A permutation-invariant LLM.
arXiv:2505.15433, 2025.

\bibitem{he2025order}
Q. He, Q. He, J. Liang, Y. Xiao, W. Zhou, Z. Sun, and F. Yu.
Order doesn't matter, but reasoning does: Training LLMs with order-centric augmentation.
arXiv:2502.19907, 2025.

\bibitem{saparov2023prontoqa}
A. Saparov and H. He.
Language models are greedy reasoners: A systematic formal analysis of chain-of-thought.
In \emph{Proceedings of ICLR}, 2023. arXiv:2210.01240.

\bibitem{merrill2024cot}
W. Merrill and A. Sabharwal.
The expressive power of transformers with chain of thought.
In \emph{Proceedings of ICLR}, 2024. arXiv:2310.07923.

\bibitem{barcelo2020gnn}
P. Barcel\'o, E. V. Kostylev, M. Monet, J. P\'erez, J. Reutter, and J. P. Silva.
The logical expressiveness of graph neural networks.
In \emph{Proceedings of ICLR}, 2020.

\bibitem{katsumata2014parametric}
S. Katsumata.
Parametric effect monads and semantics of effect systems.
In \emph{Proceedings of POPL}, pages 633--645. ACM, 2014.

\bibitem{dowling1984horn}
W. F. Dowling and J. H. Gallier.
Linear-time algorithms for testing the satisfiability of propositional Horn formulae.
\emph{Journal of Logic Programming}, 1(3):267--284, 1984.

\bibitem{nerode1958}
A. Nerode.
Linear automaton transformations.
\emph{Proceedings of the American Mathematical Society}, 9(4):541--544, 1958.

\bibitem{angluin1987learning}
D. Angluin.
Learning regular sets from queries and counterexamples.
\emph{Information and Computation}, 75(2):87--106, 1987.
\end{thebibliography}
\end{document}
